% ============================================================
%  SECCIÓN: Benchmark de EDICIÓN por instrucciones
%  Formato IEEE (IEEEtran) — español
%  Complementa benchmark_section.tex (generación). Implementación:
%  hybrid_music_engine/new_metrics/edit_benchmark.py
% ============================================================

\section{Evaluación de la Edición por Instrucciones}
\label{sec:edit}

La comparativa de generación (Sección~\ref{sec:benchmark}) justifica la
elección del modelo base, pero \emph{no} valida la tarea objetivo del sistema:
dada una pista \textbf{fuente} y una \textbf{instrucción} en lenguaje natural,
producir un audio que aplique la operación pedida \emph{preservando} el
contenido no editado de la fuente. Un buen CLAP-score en generación libre no
garantiza que se conserven identidad, instrumentación, estructura o mezcla.
Esta sección define un protocolo de evaluación específico para edición.

\subsection{Protocolo y datos}
\label{sec:edit:datos}

La unidad de evaluación es la \textbf{terna}
$\{x_{\text{src}}, \text{instr}, x_{\text{tgt}}\}$: una pista fuente, una
instrucción (p.\,ej.\ \textit{``añade una batería suave''}, \textit{``quita la
voz''}, \textit{``haz el tempo más lento''}) y un audio objetivo que
representa el resultado deseado. El editor produce una salida
$\hat{x}$ que se compara contra fuente, instrucción y objetivo.

\subsection{Líneas base}
\label{sec:edit:baselines}

Para no confundir generación con edición, se comparan al menos cuatro
referencias:
\begin{enumerate}
  \item \textbf{Reconstrucción sin edición} ($\hat{x} = x_{\text{src}}$):
        preserva el 100\,\% del contenido pero no aplica ninguna operación;
        ancla inferior de \emph{éxito de operación}.
  \item \textbf{MusicGen sólo texto}: condicionado únicamente por la
        instrucción, ignora la fuente; alto en alineación texto-audio pero
        bajo en preservación.
  \item \textbf{MusicGen con \textit{audio prompt} sin módulo de fusión}:
        usa la fuente como contexto pero sin el mecanismo de fusión aprendido.
  \item \textbf{Instruct-MusicGen} (referencia más cercana) o el adaptador
        propuesto con \texttt{AudioFusionModule}~+~LoRA.
\end{enumerate}

\subsection{Métricas}
\label{sec:edit:metricas}

Se reportan cuatro métricas automáticas y una subjetiva
(Tabla~\ref{tab:edit-metricas}):

\begin{itemize}
  \item \textbf{FAD$\to$objetivo} ($\downarrow$): distancia de Fréchet entre
        las distribuciones de embeddings de $\{\hat{x}\}$ y $\{x_{\text{tgt}}\}$.
        Cercanía al resultado deseado.
  \item \textbf{CLAP-instrucción} ($\uparrow$): coseno entre el embedding de
        $\hat{x}$ y el de la instrucción. Mide si la salida refleja lo pedido.
  \item \textbf{Preservación de contenido} ($\uparrow$): similitud coseno
        $x_{\text{src}}\!\leftrightarrow\!\hat{x}$ \emph{en lo no editado}.
        En su variante por \textit{stems} promedia la similitud sólo sobre las
        capas que la operación no debía alterar.
  \item \textbf{Éxito de operación} ($\uparrow$): fracción del cambio
        $x_{\text{src}}\!\to\!x_{\text{tgt}}$ lograda por $\hat{x}$ sobre el
        atributo objetivo (Ec.~\ref{eq:opsuccess}), recortada a $[0,1]$.
  \item \textbf{Preferencia humana} ($\uparrow$): escucha controlada tipo
        \textit{MUSHRA}~\cite{series2014method}; juez externo, no automatizable.
\end{itemize}

El éxito de operación se define como
\begin{equation}
  s \;=\; \mathrm{clip}\!\left(
  \frac{a(\hat{x}) - a(x_{\text{src}})}{a(x_{\text{tgt}}) - a(x_{\text{src}})},
  \;0,\;1\right),
  \label{eq:opsuccess}
\end{equation}
donde $a(\cdot)$ mide el atributo objetivo de la operación (energía de batería
para \textit{``añade batería''}, probabilidad de voz para \textit{``quita la
voz''}, etc.). La reconstrucción sin edición da $s=0$ (no avanza hacia el
objetivo) y una salida igual al objetivo da $s=1$.

\begin{table}[t]
  \caption{Métricas del benchmark de edición y su dirección}
  \label{tab:edit-metricas}
  \centering
  \begin{tabular}{lcl}
    \hline
    \textbf{Métrica} & \textbf{Dir.} & \textbf{Compara contra} \\
    \hline
    FAD$\to$objetivo        & $\downarrow$ & objetivo \\
    CLAP-instrucción        & $\uparrow$   & instrucción \\
    Preservación            & $\uparrow$   & fuente (no editado) \\
    Éxito de operación      & $\uparrow$   & fuente $\to$ objetivo \\
    Preferencia humana      & $\uparrow$   & juez humano \\
    \hline
  \end{tabular}
\end{table}

\subsection{Lectura conjunta}
\label{sec:edit:lectura}

Las métricas se interpretan en conjunto porque exponen una tensión que la
generación libre oculta: la reconstrucción maximiza la preservación pero
fracasa en el éxito de operación, mientras que un modelo sólo-texto puede
puntuar alto en CLAP-instrucción a costa de no preservar la fuente. Un editor
competente debe equilibrar \emph{preservar lo no editado} y \emph{aplicar la
operación}, acercándose al objetivo (FAD bajo) sin sacrificar la identidad de
la pista. El ranking global promedia los rangos de las cuatro métricas
automáticas; la preferencia humana se reporta por separado.

% ============================================================
%  Implementación: hybrid_music_engine/new_metrics/edit_benchmark.py
%  (reutiliza frechet_distance/compute_statistics de fad_score y
%   cosine_similarity de clap_score; sin reimplementar la matemática)
% ============================================================
